\chapter{ANA (Antinuclear Antibody)}

\section*{What Is This Test?}
The antinuclear antibody (ANA) test detects autoantibodies directed against components of the cell nucleus, including DNA, histones, centromeres, and other nuclear and nucleolar antigens. ANA testing is a cornerstone screening test for systemic autoimmune rheumatic diseases, particularly systemic lupus erythematosus (SLE). The test uses indirect immunofluorescence (IIF) on HEp-2 human epithelial cells, which express a broad range of nuclear antigens and provide information about both the titer (concentration) and the immunofluorescence pattern (speckled, homogeneous, nucleolar, centromere, peripheral). The pattern helps direct subsequent testing for specific autoantibodies. ANA testing is also available by enzyme immunoassay (EIA) and multiplex bead-based assays, though IIF on HEp-2 cells remains the gold standard recommended by the American College of Rheumatology \cite{agmonlevin2014international}.

\section*{Alternative Names}
ANA; Antinuclear Antibody; FANA (Fluorescent Antinuclear Antibody); ANA Screen; Antinuclear Factor (ANF); HEp-2 ANA

\section*{Principle}
The gold standard method is indirect immunofluorescence (IIF) on HEp-2 human laryngeal carcinoma cells. Patient serum is serially diluted (typically starting at 1:40 or 1:80) and incubated with HEp-2 cell substrate on glass slides. After washing, fluorescein-conjugated anti-human IgG antibody is added. The slide is examined under a fluorescence microscope by a trained technologist who reports the highest dilution at which nuclear fluorescence is still visible (titer) and the pattern of staining. Common patterns include: homogeneous (diffuse staining of entire nucleus, associated with anti-dsDNA and anti-histone antibodies), speckled (fine or coarse granular staining, associated with anti-SSA/Ro, anti-SSB/La, anti-Sm, anti-RNP), nucleolar (clumpy staining of nucleoli, associated with anti-Scl-70 in systemic sclerosis), centromere (discrete speckled pattern representing centromere proteins, associated with limited cutaneous systemic sclerosis/CREST syndrome), and peripheral/rim (staining of nuclear membrane, associated with anti-dsDNA in SLE). Solid-phase enzyme immunoassays (EIAs) and multiplex bead assays (e.g., BioPlex 2200) screen for a panel of specific autoantibodies simultaneously but do not provide pattern information.

\section*{History}
The lupus erythematosus (LE) cell phenomenon was first described by Malcolm Hargraves in 1948 \cite{hargraves1948le}, representing the earliest recognition of antinuclear antibodies. George Friou developed the indirect immunofluorescence ANA test in 1957 using rodent tissue sections as substrate \cite{friou1958interaction}. The introduction of HEp-2 human cell line substrate in the 1970s dramatically improved sensitivity, as these cells express a wider range of nuclear antigens and have larger nuclei with better visualization. The 1980s saw the identification and characterization of specific nuclear antigens (Sm, RNP, SSA/Ro, SSB/La, Scl-70, Jo-1) using immunodiffusion and immunoprecipitation. The 1982 and 1997 American College of Rheumatology (ACR) classification criteria for SLE incorporated ANA testing. The 2019 ACR/EULAR classification criteria for SLE require a positive ANA at a titer of at least 1:80 on HEp-2 cells (or an equivalent positive test) as an entry criterion for classification \cite{aringer2019eular}. The 2022 International Consensus on ANA Patterns (ICAP) standardized the nomenclature for 30 distinct HEp-2 immunofluorescence patterns \cite{damoiseaux2019clinical}.

\section*{Why This Test Is Ordered}
\begin{itemize}
  \item Screening for systemic lupus erythematosus (SLE) in patients with suggestive clinical features including malar rash, discoid rash, photosensitivity, oral ulcers, arthritis, serositis, renal disorder, neurologic disorder, hematologic disorder, or immunologic disorder
  \item Evaluation of unexplained fever, arthralgias, myalgias, Raynaud phenomenon, or constitutional symptoms suspicious for autoimmune disease
  \item Diagnosis of systemic sclerosis (scleroderma) in patients with skin thickening, Raynaud phenomenon, or interstitial lung disease
  \item Evaluation of Sjogren syndrome in patients with sicca symptoms (dry eyes and dry mouth)
  \item Diagnosis of mixed connective tissue disease (MCTD) in patients with overlapping features of SLE, systemic sclerosis, and polymyositis
  \item Diagnosis of drug-induced lupus in patients taking procainamide, hydralazine, isoniazid, minocycline, or anti-TNF agents
  \item Monitoring of autoimmune disease activity (ANA titer trends are not well-correlated with disease activity, but specific autoantibodies like anti-dsDNA are)
\end{itemize}

\section*{Is This a Primary or Secondary Test}
This is a primary screening test for systemic autoimmune rheumatic diseases. A positive ANA serves as the entry criterion for SLE classification and directs further specific autoantibody testing.

\section*{How This Test Is Performed}
For most patients, a health care professional will take a blood sample from a vein in your arm using a small needle. After the needle is inserted, a small amount of blood will be collected into a test tube or vial. You may feel a little sting when the needle goes in or out. This usually takes less than five minutes.

\section*{What Preparation Is Needed}
No special preparation is required. Fasting is not necessary. Patients should inform their healthcare provider about all medications, as certain drugs (procainamide, hydralazine, isoniazid, minocycline, anti-TNF agents, quinidine) can induce a positive ANA.

\section*{Contraindications and Precautions}
No absolute contraindications. Specimen should be collected in a serum separator tube (SST, gold-top) or red-top tube. Hemolyzed, lipemic, or icteric specimens may interfere with immunofluorescence interpretation. Specimens should be processed within 4 hours of collection or refrigerated at 2--8 degrees C for up to 48 hours. ANA positivity increases with age --- up to 15--20\% of healthy individuals over age 65 may have a low-titer positive ANA without autoimmune disease \cite{tan1997range}. Pregnancy, infections (EBV, hepatitis C, HIV, tuberculosis), and malignancies can cause false-positive ANA results. Corticosteroids and immunosuppressive medications may lower ANA titers.

\section*{Risks}
Minimal risk. Slight pain or bruising at the venipuncture site. Rare: hematoma, infection, vasovagal syncope.

\section*{What Happens After the Test}
Results are typically available within 1 to 3 days. A positive ANA result will be reported as a titer (e.g., 1:160) and a pattern (e.g., speckled). Further testing for specific autoantibodies is guided by the ANA pattern and clinical presentation. Negative ANA results generally rule out SLE (negative predictive value >95\% for titers <1:80) but do not exclude other autoimmune diseases. Your healthcare provider will interpret the results in the context of your symptoms, physical examination, and other laboratory findings.

\section*{Understanding Results}

\subsection*{Below Normal (Negative)}
\begin{itemize}
  \item A negative ANA (titer <1:40 or <1:80 depending on laboratory) in a symptomatic patient makes SLE unlikely but does not definitively exclude it --- approximately 2--5\% of SLE patients are ANA-negative, often those with anti-SSA/Ro antibodies or renal-limited disease
  \item ANA-negative lupus typically presents with prominent cutaneous manifestations and is associated with anti-SSA/Ro antibodies
  \item Drug-induced lupus (anti-histone antibodies) may be detected by specific assays even when ANA by IIF is negative
  \item Early autoimmune disease may be ANA-negative before seroconversion
  \item Immunosuppressive therapy (corticosteroids, cyclophosphamide, mycophenolate) may suppress ANA production resulting in a false-negative result
  \item Patients with convincing clinical features but negative ANA should be tested for specific autoantibodies including anti-dsDNA, anti-SSA/Ro, anti-Sm, and anti-RNP
\end{itemize}

\subsection*{Above Normal (Positive)}
\begin{itemize}
  \item Low titer ANA (1:40 or 1:80): Found in up to 20--30\% of healthy adults, particularly elderly individuals and women. Clinical significance is low in the absence of symptoms
  \item Moderate titer ANA (1:160 or 1:320): More likely to be clinically significant. Found in SLE, Sjogren syndrome, systemic sclerosis, mixed connective tissue disease, and autoimmune hepatitis
  \item High titer ANA ($\geq$1:640): Strongly associated with autoimmune disease. Very high titers ($\geq$1:1280) are characteristic of active SLE and mixed connective tissue disease
  \item Homogeneous pattern: Associated with antibodies to dsDNA, histones, and nucleosomes. Seen in SLE (anti-dsDNA), drug-induced lupus (anti-histone), and idiopathic SLE
  \item Speckled pattern: The most common pattern. Associated with antibodies to extractable nuclear antigens (ENA): anti-SSA/Ro and anti-SSB/La (Sjogren syndrome), anti-Sm and anti-RNP (SLE and MCTD), anti-Scl-70 (systemic sclerosis)
  \item Nucleolar pattern: Associated with anti-Scl-70 (diffuse systemic sclerosis), anti-PM-Scl (polymyositis-scleroderma overlap), and anti-fibrillarin/U3-RNP (systemic sclerosis)
  \item Centromere pattern: Discrete speckles (approximately 46 per cell). Highly specific for limited cutaneous systemic sclerosis (CREST syndrome --- Calcinosis, Raynaud phenomenon, Esophageal dysmotility, Sclerodactyly, Telangiectasias)
  \item Peripheral/rim pattern: Associated with anti-dsDNA antibodies and active SLE, particularly with renal involvement (lupus nephritis)
  \item False-positive ANA can occur in chronic infections (hepatitis C, HIV, EBV, tuberculosis, infective endocarditis), malignancies (hematologic and solid tumors), and with certain medications (procainamide, hydralazine, minocycline, anti-TNF biologics)
\end{itemize}

\section*{Normal Results}
\begin{itemize}
  \item \textbf{Negative:} Titer <1:40 or <1:80 (laboratory-dependent cutoff)
  \item Healthy adults (age 18--65): Negative ANA in 80--85\% of the population
  \item Healthy elderly (age >65): Negative ANA in 80--85\%; low-titer positive ANA in up to 15--20\% without clinical significance
  \item \textbf{Borderline/Equivocal:} Titer 1:40 to 1:80 (may be reported as negative by many laboratories)
  \item ANA results should always be interpreted in conjunction with the immunofluorescence pattern and specific autoantibody testing
\end{itemize}

\section*{Follow-up Tests}
\begin{itemize}
  \item \textbf{Anti-dsDNA (double-stranded DNA) antibody:} Highly specific for SLE. Rising titers correlate with disease activity and lupus nephritis flares. Measured by Farr assay (radioimmunoassay), Crithidia luciliae IIF, or ELISA
  \item \textbf{Extractable Nuclear Antigen (ENA) panel:} Includes anti-Sm (specific for SLE), anti-RNP (associated with MCTD), anti-SSA/Ro and anti-SSB/La (Sjogren syndrome, neonatal lupus, subacute cutaneous lupus), anti-Scl-70 (diffuse systemic sclerosis), and anti-Jo-1 (polymyositis)
  \item \textbf{Anti-histone antibodies:} Associated with drug-induced lupus (>95\% sensitivity)
  \item \textbf{Complement levels (C3 and C4):} Low complement levels in active SLE reflect immune complex consumption and correlate with lupus nephritis activity
  \item \textbf{Rheumatoid factor and anti-CCP antibody:} If joint symptoms are predominant, to evaluate for rheumatoid arthritis
  \item \textbf{ANCA (antineutrophil cytoplasmic antibodies):} If vasculitis is suspected
  \item \textbf{Complete blood count with differential:} To assess for anemia, leukopenia, lymphopenia, or thrombocytopenia associated with SLE
  \item \textbf{Urinalysis and spot urine protein-to-creatinine ratio:} To screen for lupus nephritis
\end{itemize}

\section*{References}
\bibliographystyle{unsrtnat}
\bibliography{references}

\clearpage
